\documentclass[a4paper,oneside,12pt]{extreport}

\usepackage{mmap}
\usepackage[T2A]{fontenc}
\usepackage[utf8]{inputenc}
\usepackage[english,russian]{babel}

% Текст отчёта следует печатать, соблюдая следующие размеры полей:
% левое — 30 мм, правое — 15 мм, верхнее и нижнее — 20 мм.
\usepackage[left=20mm, right=15mm, top=15mm, bottom=15mm]{geometry}

% \setlength{\parindent}{1.25cm} % Абзацный отступ

\usepackage{setspace}
%\onehalfspacing % Полуторный интервал

\frenchspacing % Равномерные пробелы
\usepackage{indentfirst} % Красная строка

\usepackage{microtype}
\sloppy

\usepackage{titlesec}
\titlespacing*{\chapter}{0pt}{-30pt}{8pt}
\titlespacing*{\section}{\parindent}{*4}{*4}
\titlespacing*{\subsection}{\parindent}{*4}{*4}
\titleformat{\chapter}{\LARGE\bfseries}{\thechapter}{20pt}{\LARGE\bfseries}
\titleformat{\section}{\Large\bfseries}{\thesection}{40pt}{\Large\bfseries}

\usepackage{graphicx}
\usepackage{caption}

\usepackage[unicode,pdftex]{hyperref}
\hypersetup{hidelinks}

%% title begin
\usepackage{wrapfig}

\makeatletter
	\def\vhrulefill#1{\leavevmode\leaders\hrule\@height#1\hfill \kern\z@}
\makeatother
%% title end

%% begin code
\usepackage{listings}
\usepackage{xcolor}

\lstset{
	basicstyle=\footnotesize\ttfamily,
	breakatwhitespace=true,
	breaklines=true,
	commentstyle=\color{gray},
	frame=single,
	keywordstyle=\color{blue},
	showstringspaces=false,
	stringstyle=\color{red},
	tabsize=8
}

\lstdefinestyle{lispinline}{
	frame=none,
	language=Lisp
}

\newcommand{\code}[1]{\texttt{#1}}
%% end code

%% begin theorem
\usepackage{amsthm}

\makeatletter
\newtheoremstyle{indented}
	{}% measure of space to leave above the theorem
	{}% measure of space to leave below the theorem
	{}% name of font to use in the body of the theorem
	{\parindent}% measure of space to indent
	{\bfseries}% name of head font
	{.}% punctuation between head and body
	{ }% space after theorem head; " " = normal interword space
	{}% header specification (empty for default)
\makeatother

\theoremstyle{indented}

\newtheorem{definition}{Определение}[section]
\newtheorem{example}{Пример}[section]
\newtheorem{theorem}{Теорема}[section]
\newtheorem{task}{Задание}

\makeatletter
\DeclareRobustCommand\bfseriesitshape{%
	\not@math@alphabet\itshapebfseries\relax
	\fontseries\bfdefault
	\fontshape\itdefault
	\selectfont
}
\makeatother

\DeclareTextFontCommand{\textbfit}{\bfseriesitshape}
\DeclareTextFontCommand{\define}{\bfseriesitshape}
%% end theorem

%% begin columns
\usepackage{etoolbox,refcount}
\usepackage{multicol}

\newcounter{countitems}
\newcounter{nextitemizecount}
\newcommand{\setupcountitems}{%
	\stepcounter{nextitemizecount}%
	\setcounter{countitems}{0}%
	\preto\item{\stepcounter{countitems}}%
}
\makeatletter
\newcommand{\computecountitems}{%
	\edef\@currentlabel{\number\c@countitems}%
	\label{countitems@\number\numexpr\value{nextitemizecount}-1\relax}%
}
\newcommand{\nextitemizecount}{%
	\getrefnumber{countitems@\number\c@nextitemizecount}%
}
\newcommand{\previtemizecount}{%
	\getrefnumber{countitems@\number\numexpr\value{nextitemizecount}-1\relax}%
}
\makeatother
\newenvironment{AutoMultiColItemize}{%
	\ifnumcomp{\nextitemizecount}{>}{3}{\begin{multicols}{2}}{}%
		\setupcountitems\begin{itemize}}%
		{\end{itemize}%
		\unskip\computecountitems\ifnumcomp{\previtemizecount}{>}{3}{\end{multicols}}{}}
\makeatother
\newenvironment{AutoMultiColEnumerate}{%
	\ifnumcomp{\nextitemizecount}{>}{3}{\begin{multicols}{2}}{}%
		\setupcountitems\begin{enumerate}}%
		{\end{enumerate}%
		\unskip\computecountitems\ifnumcomp{\previtemizecount}{>}{3}{\end{multicols}}{}}
%% end columns

\usepackage{array}
\usepackage{float}
\restylefloat{table}

\usepackage{amsmath}
\usepackage{amssymb}

\begin{document}

\begin{titlepage}
	{\large % 14pt instead of 12pt
	\onehalfspacing
	\centering

	\begin{wrapfigure}[7]{l}{0.14\linewidth}
		\vspace{3mm}
		\hspace{-10mm}
		\includegraphics[width=0.93\linewidth]{inc/img/bmstu-logo}
	\end{wrapfigure}
	{\singlespacing \footnotesize \bfseries Министерство науки и высшего образования Российской Федерации\\Федеральное государственное бюджетное образовательное учреждение\\высшего образования\\<<Московский государственный технический университет\\имени Н.~Э.~Баумана\\ (национальный исследовательский университет)>>\\(МГТУ им. Н.~Э.~Баумана)\\}

	\vspace{-2.2mm}
	\vhrulefill{0.9mm}\\
	\vspace{-7.5mm}
	\vhrulefill{0.2mm}\\
	\vspace{2mm}

	{\doublespacing \small \raggedright ФАКУЛЬТЕТ \hspace{25mm} «Информатика и системы управления»\\
	КАФЕДРА \hspace{5mm} «Программное обеспечение ЭВМ и информационные технологии»\\}

	\vspace{30mm}

	\textbf{ОТЧЁТ}\\
	По лабораторной работе №6\\
	По курсу: «Функциональное и логическое программирование»\\
	%Тема: «Тема работы»\\

	\vspace{60mm}

	\hspace{70mm} Студент:       \hfill Керимов~А.~Ш.\\
	\hspace{70mm} Группа:        \hfill ИУ7-64Б\\
	\hspace{70mm} Преподаватели: \hfill Толпинская~Н.~Б.,\\
	                             \hfill Строганов~Ю.~В.\\

	\vfill

	Москва\\
	\the\year\\}
\end{titlepage}

\setcounter{page}{2}


\begin{task}
	Написать предикат \code{set-equal}, который возвращает \code{t}, если два его множество-аргумента содержат одни и те же элементы, порядок которых не имеет значения.

	\begin{lstlisting}[language=Lisp, gobble=16]
		(defun set-equal (a b)
		  (null (set-exclusive-or a b)))
	\end{lstlisting}

	\begin{lstlisting}[language=Lisp, gobble=16]
		(defun set-equal (a b)
		  (if (= (list-length a) (list-length b))
		      (reduce
		        (lambda (acc cur)
		          (if acc (member cur b) nil))
		        a
		        :initial-value t)
		      nil))
	\end{lstlisting}

	\begin{lstlisting}[language=Lisp, gobble=16]
		(defun set-equal (a b)
		  (defun set-equal-process (a)
		    (cond ((null a) t)
		          ((member (car a) b) (set-equal-process (cdr a)))
		          (t nil)))
		  (if (= (list-length a) (list-length b))
		      (set-equal-process a)
		      nil))
	\end{lstlisting}
\end{task}

\begin{task}
	Напишите необходимые функции, которые обрабатывают таблицу из точечных пар: \code{(страна . столица)}, и возвращают по стране — столицу, а по столице — страну.

	\begin{lstlisting}[language=Lisp, gobble=16]
		(defun get-by-template (getter setter dict value)
		  (reduce
		    (lambda (acc cur)
		      (cond (acc acc)
		            ((equal value (funcall getter cur)) (funcall setter cur))
		            (t nil)))
		    dict
		   :initial-value nil))

		(defun get-capital-by-country (dict country)
		  (get-by-template #'car #'cdr dict country))

		(defun get-country-by-capital (dict capital)
		  (get-by-template #'cdr #'car dict capital))
	\end{lstlisting}

	\begin{lstlisting}[language=Lisp, gobble=16]
		(defun get-by-template (getter setter dict value)
		  (cond ((null dict) nil)
		        ((equal (funcall getter dict) value) (funcall setter dict))
		        (t (get-by-template getter setter (cdr dict) value))))

		(defun get-capital-by-country (dict country)
		  (get-by-template #'caar #'cdar dict country))

		(defun get-country-by-capital (dict capital)
		  (get-by-template #'cdar #'caar dict capital))
	\end{lstlisting}
\end{task}

\begin{task}
	Напишите функцию, которая умножает на заданное число-аргумент все числа из заданного списка-аргумента, когда а) все элементы списка — числа, б)  элементы списка — любые объекты.

	\begin{lstlisting}[language=Lisp, gobble=16]
		(defun mult-a (lst number)
		  (mapcar
		    (lambda (item)
		      (* item number))
		  lst))

		(defun mult-b (lst number)
		  (mapcar
		    (lambda (item)
		      (if (numberp item)
		          (* item number)
		          item))
		  lst))
	\end{lstlisting}

	\begin{lstlisting}[language=Lisp, gobble=16]
		(defun mult-a (lst number)
		  (defun mult-a-process (lst acc)
		    (if lst
		        (mult-a-process
		          (cdr lst)
		          (append acc
		            (list (* (car lst) number))))
		        acc))
		  (mult-a-process lst nil))

		(defun mult-b (lst number)
		  (defun mult-b-process (lst acc)
		    (if lst
		        (mult-b-process
		          (cdr lst)
		          (append acc
		            (if (numberp (car lst))
		                (list (* number (car lst)))
		                (list (car lst)))))
		        acc))
		  (mult-b-process lst nil))
	\end{lstlisting}
\end{task}

\begin{task}
	Напишите функцию, которая уменьшает на 10 все числа из списка аргумента этой функции.

	\begin{lstlisting}[language=Lisp, gobble=16]
		(defun minus (lst number)
		  (mapcar
		    (lambda (item)
		      (- item number))
		    lst))

		(defun minus-ten (lst)
		  (minus lst 10))
	\end{lstlisting}

	\begin{lstlisting}[language=Lisp, gobble=16]
		(defun minus (lst number)
		  (defun minus-process (lst acc)
		    (if lst
		        (minus-process
		          (cdr lst)
		          (append acc
		            (list (- (car lst) number))))
		        acc))
		  (minus-process lst nil))

		(defun minus-ten (lst)
		  (minus lst 10))
	\end{lstlisting}
\end{task}

\begin{task}
	Написать функцию, которая возвращает первый аргумент списка-аргумента, который сам является непустым списком.

	\begin{lstlisting}[language=Lisp, gobble=16]
		(defun get-first-sublist (lst)
		  (reduce
		    (lambda (acc cur)
		      (cond (acc acc)
		            ((listp cur) cur)
		            (t nil)))
		    lst
		    :initial-value nil))
	\end{lstlisting}

	\begin{lstlisting}[language=Lisp, gobble=16]
		(defun not-empty-list (lst)
		  (and (listp lst) lst))

		(defun get-first-sublist (lst)
		  (cond ((null lst) nil)
		        ((not-empty-list (car lst)) (car lst))
		        (t (get-first-sublist (cdr lst)))))
	\end{lstlisting}
\end{task}

\begin{task}
	Написать функцию, которая выбирает из заданного списка только те числа, которые больше 1 и меньше 10.
	(Вариант: между двумя заданными границами)

	\begin{lstlisting}[language=Lisp, gobble=16]
		(defun between (lst a b)
		  (reduce
		    (lambda (acc cur)
		      (if (and (< a cur) (< cur b))
		          (append acc (list cur))
		          acc))
		    lst
		    :initial-value nil))

		(defun between-1-10 (lst)
		  (between lst 1 10))
	\end{lstlisting}

	\begin{lstlisting}[language=Lisp, gobble=16]
		(defun between (lst a b)
		  (defun is-between (number)
		    (and (< a number) (< number b)))
		  (defun between-process (lst acc)
		    (if lst
		        (between-process
		          (cdr lst)
		          (if (is-between (car lst))
		              (append acc (list (car lst)))
		              acc))
		        acc))
		  (between-process lst nil))

		(defun between-1-10 (lst)
		  (between lst 1 10))
	\end{lstlisting}
\end{task}

\begin{task}
	Написать функцию, вычисляющую декартово произведение двух своих списков-аргументов.
	(Напомним, что $A\times B$ — это множество всевозможных пар $(a\;b)$, где $a\in A$, $b\in B$)

	\begin{lstlisting}[language=Lisp, gobble=16]
		(defun set-mult (a b)
		  (apply
		    #'append
		    (mapcar
		      (lambda (item-a)
		        (mapcar
		          (lambda (item-b)
		            (list item-a item-b))
		          b))
		      a)))
	\end{lstlisting}

	\begin{lstlisting}[language=Lisp, gobble=16]
		(defun set-mult (a b)
		  (defun set-mul-process (x y)
		    (cond ((null x) nil)
		          ((null y) (set-mul-process (cdr x) b))
		          (t (cons (list (car x) (car y))
		                   (set-mul-process x (cdr y))))))
		  (set-mul-process a b))
	\end{lstlisting}
\end{task}

\begin{task}
	Почему так реализовано \code{reduce}, в чем причина?

	\begin{lstlisting}[language=Lisp, gobble=16]
		(reduce #'+ 0)  ; 0
		(reduce #'+ ()) ; 0
	\end{lstlisting}
\end{task}

\end{document}
